\documentclass[10pt,a4paper]{amsart}
\usepackage{amsmath,amssymb,amsthm,graphicx}
% Retained minimal preamble for the Stacks Project chapter more-algebra.tex (More on Algebra).
% The chapter file pinned-original-source.tex carries only \input{preamble} on its first line, so no
% preamble text lives inside the chapter itself.  The Stacks Project uses one single shared preamble for
% every chapter; the snapshot stages it identically as sources/stacks-chow/preamble.tex and
% sources/stacks-etale-cohomology/preamble.tex.  The block below is copied mechanically, line for line,
% from sources/stacks-chow/preamble.tex lines 14-18 (the Xy-pic setup) and 154-217 (the native theorem
% declarations together with the mathematical macros).  Nothing was added, removed or redefined; the
% wrapper class and the amsmath/amssymb/amsthm/graphicx base packages are supplied by the assembler.
% Chapter: more-algebra.tex -- More on Algebra.
% For commutative diagrams we use Xy-pic
\usepackage[all]{xy}

% We use 2cell for 2-commutative diagrams.
\xyoption{2cell}

% Theorem environments.
%
\theoremstyle{plain}
\newtheorem{theorem}[subsection]{Theorem}
\newtheorem{proposition}[subsection]{Proposition}
\newtheorem{lemma}[subsection]{Lemma}

\theoremstyle{definition}
\newtheorem{definition}[subsection]{Definition}
\newtheorem{example}[subsection]{Example}
\newtheorem{exercise}[subsection]{Exercise}
\newtheorem{situation}[subsection]{Situation}

\theoremstyle{remark}
\newtheorem{remark}[subsection]{Remark}
\newtheorem{remarks}[subsection]{Remarks}

\numberwithin{equation}{subsection}

% Macros
%
\def\lim{\mathop{\mathrm{lim}}\nolimits}
\def\colim{\mathop{\mathrm{colim}}\nolimits}
\def\Spec{\mathop{\mathrm{Spec}}}
\def\Hom{\mathop{\mathrm{Hom}}\nolimits}
\def\Ext{\mathop{\mathrm{Ext}}\nolimits}
\def\SheafHom{\mathop{\mathcal{H}\!\mathit{om}}\nolimits}
\def\SheafExt{\mathop{\mathcal{E}\!\mathit{xt}}\nolimits}
\def\Sch{\mathit{Sch}}
\def\Mor{\mathop{\mathrm{Mor}}\nolimits}
\def\Ob{\mathop{\mathrm{Ob}}\nolimits}
\def\Sh{\mathop{\mathit{Sh}}\nolimits}
\def\NL{\mathop{N\!L}\nolimits}
\def\CH{\mathop{\mathrm{CH}}\nolimits}
\def\proetale{{pro\text{-}\acute{e}tale}}
\def\etale{{\acute{e}tale}}
\def\QCoh{\mathit{QCoh}}
\def\Ker{\mathop{\mathrm{Ker}}}
\def\Im{\mathop{\mathrm{Im}}}
\def\Coker{\mathop{\mathrm{Coker}}}
\def\Coim{\mathop{\mathrm{Coim}}}

% Boxtimes
%
\DeclareMathSymbol{\boxtimes}{\mathbin}{AMSa}{"02}

%
% Macros for moduli stacks/spaces
%
\def\QCohstack{\mathcal{QC}\!\mathit{oh}}
\def\Cohstack{\mathcal{C}\!\mathit{oh}}
\def\Spacesstack{\mathcal{S}\!\mathit{paces}}
\def\Quotfunctor{\mathrm{Quot}}
\def\Hilbfunctor{\mathrm{Hilb}}
\def\Curvesstack{\mathcal{C}\!\mathit{urves}}
\def\Polarizedstack{\mathcal{P}\!\mathit{olarized}}
\def\Complexesstack{\mathcal{C}\!\mathit{omplexes}}
% \Pic is the operator that assigns to X its picard group, usage \Pic(X)
% \Picardstack_{X/B} denotes the Picard stack of X over B
% \Picardfunctor_{X/B} denotes the Picard functor of X over B
\def\Pic{\mathop{\mathrm{Pic}}\nolimits}
\def\Picardstack{\mathcal{P}\!\mathit{ic}}
\def\Picardfunctor{\mathrm{Pic}}
\def\Deformationcategory{\mathcal{D}\!\mathit{ef}}

% --- appended 2026-10-12 (Chief): the project own definitions of the reference/slogan/history
% display environments, copied mechanically from the shared Stacks preamble lines 8-12
% (\usepackage{verbatim} plus the three \newenvironment lines, which the chapter preamble keeps
% verbatim from the project preamble). Nothing added, removed or redefined: these environments
% carry no mathematical content and the project itself typesets them as comments. Without these
% three lines a counted statement whose own text uses them fails with Environment slogan undefined.
\usepackage{verbatim}
\newenvironment{reference}{\comment}{\endcomment}
\newenvironment{slogan}{\comment}{\endcomment}
\newenvironment{history}{\comment}{\endcomment}
\usepackage{cleveref}
\crefname{originallemma}{original lemma}{original lemmas}
\Crefname{originallemma}{Original lemma}{Original lemmas}
\makeatletter
\newlabel{lemma-Rlim-pro-equal}{{23405}{source}{}{sourcepointer-1}{}}
\newlabel{lemma-Rlim-pro-equal@cref}{{[originallemma][0][]23405}{[1][0][]source}{}{}{}}
\makeatother
\crefname{originallemma}{original lemma}{original lemmas}
\Crefname{originallemma}{Original lemma}{Original lemmas}
\makeatletter
\newlabel{lemma-Rlim-zero-of-direct-sums}{{23527}{source}{}{sourcepointer-2}{}}
\newlabel{lemma-Rlim-zero-of-direct-sums@cref}{{[originallemma][0][]23527}{[1][0][]source}{}{}{}}
\makeatother
\crefname{originallemma}{original lemma}{original lemmas}
\Crefname{originallemma}{Original lemma}{Original lemmas}
\makeatletter
\newlabel{lemma-derived-completion-koszul}{{26388}{source}{}{sourcepointer-3}{}}
\newlabel{lemma-derived-completion-koszul@cref}{{[originallemma][0][]26388}{[1][0][]source}{}{}{}}
\makeatother
\crefname{originallemma}{original lemma}{original lemmas}
\Crefname{originallemma}{Original lemma}{Original lemmas}
\makeatletter
\newlabel{lemma-lift-universally}{{26782}{source}{}{sourcepointer-4}{}}
\newlabel{lemma-lift-universally@cref}{{[originallemma][0][]26782}{[1][0][]source}{}{}{}}
\makeatother
\crefname{originallemma}{original lemma}{original lemmas}
\Crefname{originallemma}{Original lemma}{Original lemmas}
\makeatletter
\newlabel{lemma-tensor-Rlim-exact}{{23915}{source}{}{sourcepointer-5}{}}
\newlabel{lemma-tensor-Rlim-exact@cref}{{[originallemma][0][]23915}{[1][0][]source}{}{}{}}
\makeatother
\title{Naive derived completion and derived completion agree exactly when the f-power torsion is bounded}
\author{The Stacks Project Authors}
\date{stacks-project:more-algebra.tex:local-snapshot-20261012; GNU FDL 1.2 or later}
\begin{document}
\maketitle
\setcounter{subsection}{1}
\noindent\textit{Editorial note.} All delivered mathematical text is the author's own, copied from the pinned Stacks Project chapter file more-algebra.tex; no statement or proof step was rewritten, reordered or removed.\par\medskip
\noindent\textit{Editorial note.} The chapter file pinned-original-source.tex carries only \textbackslash{}input\{preamble\} on its first line, so no preamble text lives inside the chapter. The retained file sources/stacks-more-algebra/retained-preamble.tex is a mechanical, line-for-line copy of the project's single shared preamble as staged in the snapshot; nothing was added, removed or redefined and no package was added.\par\medskip
\noindent\textit{Editorial note.} The source URL is the source-level Stacks Project URL, as required when more than one card is produced from one source; the per-claim tag is carried by claim\_key.\par\medskip
\noindent\textit{Editorial note.} The counted statement span displays the derived completion via the two-term system (A -> f\textasciicircum{}n -> A); that displayed material is delivered verbatim. The five references inside the counted statement and proof spans are recorded as non-load-bearing navigation pointers.\par\medskip
\noindent\textit{Allowed standard prerequisites (editorial).}\begin{itemize}
\item The lifting statement identifying the pro-object of multiplication maps with the pro-object of quotient rings when the torsion is bounded. invoked at proof lines 26825-26827. Gives the forward implication.
\item The exactness statement for tensor products with R lim and the vanishing criterion. invoked at proof lines 26829-26834. Turns the isomorphism condition into a vanishing condition for all K.
\item Vanishing of the derived limit of direct sums in the relevant range. invoked at proof lines 26836-26838. Forces the system of f-power torsion modules to be pro-zero.
\item The derived completion defined via the Koszul-type system (A -> f\textasciicircum{}n -> A). referred to at statement lines 26816-26819. Fixes the left-hand functor of the natural transformation being compared.
\end{itemize}\medskip
\ifdefined\hypertarget\hypertarget{sourcepointer-1}{}\fi
\noindent\textit{Source reference.} Original lemma 23405: sources/stacks-more-algebra/pinned-original-source.tex, source lines 23405--23405. Same-chapter navigation pointer; the statement it names is standard for the section and the argument that consumes it stays inside the counted proof span.\par\medskip
\ifdefined\hypertarget\hypertarget{sourcepointer-2}{}\fi
\noindent\textit{Source reference.} Original lemma 23527: sources/stacks-more-algebra/pinned-original-source.tex, source lines 23527--23527. Same-chapter navigation pointer; the statement it names is standard for the section and the argument that consumes it stays inside the counted proof span.\par\medskip
\ifdefined\hypertarget\hypertarget{sourcepointer-3}{}\fi
\noindent\textit{Source reference.} Original lemma 26388: sources/stacks-more-algebra/pinned-original-source.tex, source lines 26388--26388. Same-chapter navigation pointer; the statement it names is standard for the section and the argument that consumes it stays inside the counted proof span.\par\medskip
\ifdefined\hypertarget\hypertarget{sourcepointer-4}{}\fi
\noindent\textit{Source reference.} Original lemma 26782: sources/stacks-more-algebra/pinned-original-source.tex, source lines 26782--26782. Same-chapter navigation pointer; the statement it names is standard for the section and the argument that consumes it stays inside the counted proof span.\par\medskip
\ifdefined\hypertarget\hypertarget{sourcepointer-5}{}\fi
\noindent\textit{Source reference.} Original lemma 23915: sources/stacks-more-algebra/pinned-original-source.tex, source lines 23915--23915. Same-chapter navigation pointer; the statement it names is standard for the section and the argument that consumes it stays inside the counted proof span.\par\medskip
\begin{lemma}
\label{lemma-when-does-it-work}
Let $A$ be a ring and $f \in A$. Set $I = (f)$. In this situation
we have the naive derived completion
$K \mapsto K' = R\lim (K \otimes_A^\mathbf{L} A/f^nA)$ and the
derived completion
$$
K \mapsto K^\wedge = R\lim (K \otimes_A^\mathbf{L} (A \xrightarrow{f^n} A))
$$
of Lemma \ref{lemma-derived-completion-koszul}.
The natural transformation of functors $K^\wedge \to K'$
is an isomorphism if and only if the $f$-power torsion of $A$ is bounded.
\end{lemma}

\begin{proof}
If the $f$-power torsion is bounded, then the pro-objects
$\{(f^n : A \to A)\}$ and $\{A/f^nA\}$ are isomorphic by
Lemma \ref{lemma-lift-universally}.
Hence the functors are isomorphic by Lemma \ref{lemma-Rlim-pro-equal}.
Conversely, we see from Lemma \ref{lemma-tensor-Rlim-exact}
that the condition is exactly that
$$
R\lim (K \otimes_A^\mathbf{L} A[f^n])
$$
is zero for all $K \in D(A)$. Here the maps of the system $(A[f^n])$
are given by multiplication by $f$. Taking $K = A$ and
$K = \bigoplus_{i \in \mathbf{N}} A$ we see from
Lemma \ref{lemma-Rlim-zero-of-direct-sums}
this implies $(A[f^n])$ is zero as a pro-object, i.e.,
$f^{n - 1}A[f^n] = 0$ for some $n$, i.e., $A[f^{n - 1}] = A[f^n]$, i.e.,
the $f$-power torsion is bounded.
\end{proof}

\end{document}
